\documentclass{beamer}\usepackage[]{graphicx}\usepackage[]{color}
%% maxwidth is the original width if it is less than linewidth
%% otherwise use linewidth (to make sure the graphics do not exceed the margin)
\makeatletter
\def\maxwidth{ %
  \ifdim\Gin@nat@width>\linewidth
    \linewidth
  \else
    \Gin@nat@width
  \fi
}
\makeatother

\definecolor{fgcolor}{rgb}{1, 0.894, 0.769}
\newcommand{\hlnum}[1]{\textcolor[rgb]{0.824,0.412,0.118}{#1}}%
\newcommand{\hlstr}[1]{\textcolor[rgb]{1,0.894,0.71}{#1}}%
\newcommand{\hlcom}[1]{\textcolor[rgb]{0.824,0.706,0.549}{#1}}%
\newcommand{\hlopt}[1]{\textcolor[rgb]{1,0.894,0.769}{#1}}%
\newcommand{\hlstd}[1]{\textcolor[rgb]{1,0.894,0.769}{#1}}%
\newcommand{\hlkwa}[1]{\textcolor[rgb]{0.941,0.902,0.549}{#1}}%
\newcommand{\hlkwb}[1]{\textcolor[rgb]{0.804,0.776,0.451}{#1}}%
\newcommand{\hlkwc}[1]{\textcolor[rgb]{0.78,0.941,0.545}{#1}}%
\newcommand{\hlkwd}[1]{\textcolor[rgb]{1,0.78,0.769}{#1}}%
\let\hlipl\hlkwb

\usepackage{framed}
\makeatletter
\newenvironment{kframe}{%
 \def\at@end@of@kframe{}%
 \ifinner\ifhmode%
  \def\at@end@of@kframe{\end{minipage}}%
  \begin{minipage}{\columnwidth}%
 \fi\fi%
 \def\FrameCommand##1{\hskip\@totalleftmargin \hskip-\fboxsep
 \colorbox{shadecolor}{##1}\hskip-\fboxsep
     % There is no \\@totalrightmargin, so:
     \hskip-\linewidth \hskip-\@totalleftmargin \hskip\columnwidth}%
 \MakeFramed {\advance\hsize-\width
   \@totalleftmargin\z@ \linewidth\hsize
   \@setminipage}}%
 {\par\unskip\endMakeFramed%
 \at@end@of@kframe}
\makeatother

\definecolor{shadecolor}{rgb}{.97, .97, .97}
\definecolor{messagecolor}{rgb}{0, 0, 0}
\definecolor{warningcolor}{rgb}{1, 0, 1}
\definecolor{errorcolor}{rgb}{1, 0, 0}
\newenvironment{knitrout}{}{} % an empty environment to be redefined in TeX

\usepackage{alltt}
\usepackage{../371g-slides}
\title{Model Building - Part 2}
\subtitle{Lecture 15}
\author{STA 371G}
\IfFileExists{upquote.sty}{\usepackage{upquote}}{}
\begin{document}
  
  

  \frame{\maketitle}

  % Show outline at beginning of each section
  \AtBeginSection[]{ 
    \begin{frame}<beamer>
      \tableofcontents[currentsection]
    \end{frame}
  }

  %%%%%%% Slides start here %%%%%%%

  \begin{darkframes}
    
    %{slide 1}
    \begin{frame}{Predicting Baseball Player Batting Averages}
      \fontsize{10}{10}\selectfont
      \begin{center}
        \includegraphics[width=1.6in]{TyCobb.png} \\
      \end{center}

      \begin{itemize}
        \item load baseball.csv
        \item install the packages car, leaps, and corrplot if you haven't already
      \end{itemize}
      
      \lc %{smallest population}
    \end{frame}


    %{slide 2}
    \begin{frame}[fragile]{What predicts a player's batting average}
      \begin{itemize}[<+->]
        \item All of the data here came from http://seanlahman.com/baseball-archive/statistics/
        \item Some data cleaning, to calculate averages mostly, was done.
        \item We are going to explore this dataset with best subsets regression
        \end{itemize} 
    \end{frame}


    %{slide 3}
    \begin{frame}[fragile]{The 10 potential x variables}
    \fontsize{10}{10}\selectfont
          \begin{itemize}
            \item YEAR:    Year this entry calculated for 
            \item LG:      League, either NL or AL
            \item AVG:     Batting average
            \item OBP:     On base percentage
            \item SLG:     Slugging average
            \item EXP:     Years of experience
            \item PAYR:    Plate appearances per year
            \item MLAVG:   Batting average for the leauge for the year
            \item MLOBP:   On base percentage for the leaugue for the year
            \item MLSLG:   Slugging percentage for the leaugue for the year
            \item AVGcumLag1:   Player's cumulative batting average for previous years
            \item OBPcumLag1:   Player's cumulative on base percentage for previous years
            \item SLGcumLag1:   Player's cumulative slugging percentage for previous years
            \item G:       Games played (must have been at least 98)
            \item YRINDEX: Number of years since 1958
          \end{itemize}

        \lc %{Why predict batting average instead of hits?}
    \end{frame}


    %{slide 4}
    \begin{frame}[fragile]{Build model full and check for multicollinearity}
      \fontsize{8}{8}\selectfont  

\begin{knitrout}
\definecolor{shadecolor}{rgb}{0.137, 0.137, 0.137}\begin{kframe}
\begin{alltt}
\hlstd{full} \hlkwb{<-} \hlkwd{lm}\hlstd{(AVG} \hlopt{~} \hlstd{OBP} \hlopt{+} \hlstd{SLG} \hlopt{+} \hlstd{EXP} \hlopt{+} \hlstd{PAYR} \hlopt{+} \hlstd{MLAVG}
\hlopt{+} \hlstd{MLOBP} \hlopt{+} \hlstd{MLSLG} \hlopt{+} \hlstd{AVGcumLag1} \hlopt{+} \hlstd{OBPcumLag1}
\hlopt{+} \hlstd{SLGcumLag1} \hlopt{+} \hlstd{G} \hlopt{+} \hlstd{YRINDEX,} \hlkwc{data}\hlstd{=baseball)}

\hlkwd{round}\hlstd{(}\hlkwd{vif}\hlstd{(full),}\hlnum{2}\hlstd{)}
\end{alltt}
\begin{verbatim}
       OBP        SLG        EXP       PAYR      MLAVG      MLOBP 
      3.71       4.32       1.20       1.37      11.07      12.69 
     MLSLG AVGcumLag1 OBPcumLag1 SLGcumLag1          G    YRINDEX 
      7.39       2.09       3.95       3.82       1.12       2.18 
\end{verbatim}
\end{kframe}
\end{knitrout}
      
      \pause
      \fontsize{10}{10}\selectfont{Uh oh. Houston, we have a problem!}

      \lc %{What is the R2 of model full?}
    \end{frame}


    %{slide #5}
    \begin{frame}[fragile]{Look at the correlations to find the problem}
      \fontsize{8}{8}\selectfont  

      {This matrix is hard to read}
\begin{knitrout}
\definecolor{shadecolor}{rgb}{0.137, 0.137, 0.137}\begin{kframe}
\begin{alltt}
\hlstd{numericpredictors} \hlkwb{<-} \hlstd{baseball[,}\hlnum{8}\hlopt{:}\hlnum{19}\hlstd{]}
\hlstd{M} \hlkwb{<-} \hlkwd{round}\hlstd{(}\hlkwd{cor}\hlstd{(numericpredictors),}\hlnum{2}\hlstd{)} \hlcom{# calculate correlations}

\hlcom{# print M by just typing M}

\hlcom{# This table is confusing!}
\hlcom{# There is a much better way to see this using corrplot}
\hlcom{# So make the library available}
\hlkwd{library}\hlstd{(corrplot)}
\end{alltt}
\end{kframe}
\end{knitrout}
      
      \lc %{What is the correlation between OBP and SLG?}
    \end{frame}


    %{slide #6}
    \begin{frame}[fragile]{Plot the correlations to better see the problem}
      \fontsize{8}{8}\selectfont  

\begin{knitrout}
\definecolor{shadecolor}{rgb}{0.137, 0.137, 0.137}\begin{kframe}
\begin{alltt}
\hlkwd{corrplot}\hlstd{(M,} \hlkwc{method} \hlstd{=} \hlstr{"circle"}\hlstd{)} \hlcom{#plot matrix}
\end{alltt}
\end{kframe}
\input{c:/tmp/figures/unnamed-chunk-4-1.tikz}

\end{knitrout}
    \end{frame}

    %{slide 7}
    \begin{frame}[fragile]{Reduce multicollinearity by dropping variables}
      \fontsize{8}{8}\selectfont 
      \begin{itemize}
        \item The Major League averages are highly correlated with each other
        \item Let's keep just MLAVG and drop MLOBP and MLSLG
        \end{itemize} 

\begin{knitrout}
\definecolor{shadecolor}{rgb}{0.137, 0.137, 0.137}\begin{kframe}
\begin{alltt}
\hlstd{full} \hlkwb{<-} \hlkwd{lm}\hlstd{(AVG} \hlopt{~} \hlstd{OBP} \hlopt{+} \hlstd{SLG} \hlopt{+} \hlstd{EXP} \hlopt{+} \hlstd{PAYR} \hlopt{+} \hlstd{MLAVG}
\hlopt{+} \hlstd{AVGcumLag1} \hlopt{+} \hlstd{OBPcumLag1}
\hlopt{+} \hlstd{SLGcumLag1} \hlopt{+} \hlstd{G} \hlopt{+} \hlstd{YRINDEX,} \hlkwc{data}\hlstd{=baseball)}

\hlkwd{round}\hlstd{(}\hlkwd{vif}\hlstd{(full),} \hlnum{2}\hlstd{)}
\end{alltt}
\begin{verbatim}
       OBP        SLG        EXP       PAYR      MLAVG AVGcumLag1 
      3.62       4.29       1.16       1.37       1.86       2.09 
OBPcumLag1 SLGcumLag1          G    YRINDEX 
      3.92       3.79       1.12       1.85 
\end{verbatim}
\end{kframe}
\end{knitrout}

      {Much better!}
    \end{frame}


    %{slide 8}
    \begin{frame}[fragile]{Run resubsets to get a sense of the best predictors}
    \fontsize{8}{8}\selectfont  
\begin{knitrout}
\definecolor{shadecolor}{rgb}{0.137, 0.137, 0.137}\begin{kframe}
\begin{alltt}
\hlkwd{library}\hlstd{(leaps)}
\hlstd{bestsubsets} \hlkwb{<-} \hlkwd{regsubsets}\hlstd{(AVG} \hlopt{~} \hlstd{OBP} \hlopt{+} \hlstd{SLG} \hlopt{+} \hlstd{EXP} \hlopt{+} \hlstd{PAYR} \hlopt{+} \hlstd{MLAVG}
\hlopt{+} \hlstd{AVGcumLag1} \hlopt{+} \hlstd{OBPcumLag1}
\hlopt{+} \hlstd{SLGcumLag1} \hlopt{+} \hlstd{G} \hlopt{+} \hlstd{YRINDEX,} \hlkwc{data}\hlstd{=baseball)}
\end{alltt}
\end{kframe}
\end{knitrout}

      {Now let's plot and identify the important predictors}
    \end{frame}


    %{slide 9}
    \begin{frame}[fragile]{Use Adj $R^2$ to compare models}
      \fontsize{8}{8}\selectfont  
\begin{knitrout}
\definecolor{shadecolor}{rgb}{0.137, 0.137, 0.137}\begin{kframe}
\begin{alltt}
\hlkwd{plot}\hlstd{(bestsubsets,} \hlkwc{scale}\hlstd{=}\hlstr{"adjr2"}\hlstd{)}  \hlcom{# use adjusted R^2}
\end{alltt}
\end{kframe}
% Created by tikzDevice version 0.10.1 on 2017-03-09 02:37:31
% !TEX encoding = UTF-8 Unicode
\begin{tikzpicture}[x=1pt,y=1pt]
\definecolor{fillColor}{RGB}{255,255,255}
\path[use as bounding box,fill=fillColor,fill opacity=0.00] (0,0) rectangle (289.08,216.81);
\begin{scope}
\path[clip] (  0.00,  0.00) rectangle (289.08,216.81);
\definecolor{drawColor}{RGB}{254,254,254}

\path[draw=drawColor,line width= 0.4pt,line join=round,line cap=round] ( 61.20, 85.20) --
	(251.88, 85.20) --
	(251.88,143.61) --
	( 61.20,143.61) --
	( 61.20, 85.20);
\end{scope}
\begin{scope}
\path[clip] (  0.00,  0.00) rectangle (289.08,216.81);
\definecolor{drawColor}{RGB}{254,254,254}

\node[text=drawColor,rotate= 90.00,anchor=base,inner sep=0pt, outer sep=0pt, scale=  1.00] at ( 22.80,114.40) {adjr2};
\end{scope}
\begin{scope}
\path[clip] ( 61.20, 85.20) rectangle (251.88,143.61);
\definecolor{fillColor}{RGB}{230,230,230}

\path[fill=fillColor] ( 61.20, 85.20) rectangle ( 78.53, 92.50);
\definecolor{fillColor}{gray}{0.70}

\path[fill=fillColor] ( 61.20, 92.50) rectangle ( 78.53, 99.80);
\definecolor{fillColor}{gray}{0.10}

\path[fill=fillColor] ( 61.20, 99.80) rectangle ( 78.53,107.10);

\path[fill=fillColor] ( 61.20,107.10) rectangle ( 78.53,114.40);
\definecolor{fillColor}{RGB}{0,0,0}

\path[fill=fillColor] ( 61.20,114.40) rectangle ( 78.53,121.71);

\path[fill=fillColor] ( 61.20,121.71) rectangle ( 78.53,129.01);

\path[fill=fillColor] ( 61.20,129.01) rectangle ( 78.53,136.31);

\path[fill=fillColor] ( 61.20,136.31) rectangle ( 78.53,143.61);
\definecolor{fillColor}{RGB}{230,230,230}

\path[fill=fillColor] ( 78.53, 85.20) rectangle ( 95.87, 92.50);
\definecolor{fillColor}{gray}{0.70}

\path[fill=fillColor] ( 78.53, 92.50) rectangle ( 95.87, 99.80);
\definecolor{fillColor}{gray}{0.10}

\path[fill=fillColor] ( 78.53, 99.80) rectangle ( 95.87,107.10);

\path[fill=fillColor] ( 78.53,107.10) rectangle ( 95.87,114.40);
\definecolor{fillColor}{RGB}{0,0,0}

\path[fill=fillColor] ( 78.53,114.40) rectangle ( 95.87,121.71);

\path[fill=fillColor] ( 78.53,121.71) rectangle ( 95.87,129.01);

\path[fill=fillColor] ( 78.53,129.01) rectangle ( 95.87,136.31);

\path[fill=fillColor] ( 78.53,136.31) rectangle ( 95.87,143.61);
\definecolor{fillColor}{gray}{0.10}

\path[fill=fillColor] ( 95.87,107.10) rectangle (113.20,114.40);
\definecolor{fillColor}{RGB}{0,0,0}

\path[fill=fillColor] ( 95.87,114.40) rectangle (113.20,121.71);

\path[fill=fillColor] ( 95.87,121.71) rectangle (113.20,129.01);

\path[fill=fillColor] ( 95.87,129.01) rectangle (113.20,136.31);

\path[fill=fillColor] ( 95.87,136.31) rectangle (113.20,143.61);

\path[fill=fillColor] (113.20,129.01) rectangle (130.54,136.31);

\path[fill=fillColor] (113.20,136.31) rectangle (130.54,143.61);
\definecolor{fillColor}{gray}{0.70}

\path[fill=fillColor] (165.21, 92.50) rectangle (182.54, 99.80);
\definecolor{fillColor}{gray}{0.10}

\path[fill=fillColor] (165.21, 99.80) rectangle (182.54,107.10);

\path[fill=fillColor] (165.21,107.10) rectangle (182.54,114.40);
\definecolor{fillColor}{RGB}{0,0,0}

\path[fill=fillColor] (165.21,114.40) rectangle (182.54,121.71);

\path[fill=fillColor] (165.21,121.71) rectangle (182.54,129.01);

\path[fill=fillColor] (165.21,129.01) rectangle (182.54,136.31);

\path[fill=fillColor] (165.21,136.31) rectangle (182.54,143.61);
\definecolor{fillColor}{gray}{0.10}

\path[fill=fillColor] (182.54, 99.80) rectangle (199.88,107.10);

\path[fill=fillColor] (182.54,107.10) rectangle (199.88,114.40);
\definecolor{fillColor}{RGB}{0,0,0}

\path[fill=fillColor] (182.54,114.40) rectangle (199.88,121.71);

\path[fill=fillColor] (182.54,121.71) rectangle (199.88,129.01);

\path[fill=fillColor] (182.54,129.01) rectangle (199.88,136.31);

\path[fill=fillColor] (182.54,136.31) rectangle (199.88,143.61);

\path[fill=fillColor] (199.88,114.40) rectangle (217.21,121.71);

\path[fill=fillColor] (199.88,121.71) rectangle (217.21,129.01);

\path[fill=fillColor] (199.88,129.01) rectangle (217.21,136.31);

\path[fill=fillColor] (199.88,136.31) rectangle (217.21,143.61);

\path[fill=fillColor] (217.21,136.31) rectangle (234.55,143.61);

\path[fill=fillColor] (234.55,121.71) rectangle (251.88,129.01);

\path[fill=fillColor] (234.55,129.01) rectangle (251.88,136.31);

\path[fill=fillColor] (234.55,136.31) rectangle (251.88,143.61);
\end{scope}
\begin{scope}
\path[clip] (  0.00,  0.00) rectangle (289.08,216.81);
\definecolor{drawColor}{RGB}{254,254,254}

\path[draw=drawColor,line width= 0.4pt,line join=round,line cap=round] ( 69.87, 85.20) -- (243.21, 85.20);

\path[draw=drawColor,line width= 0.4pt,line join=round,line cap=round] ( 69.87, 85.20) -- ( 69.87, 79.20);

\path[draw=drawColor,line width= 0.4pt,line join=round,line cap=round] ( 87.20, 85.20) -- ( 87.20, 79.20);

\path[draw=drawColor,line width= 0.4pt,line join=round,line cap=round] (104.54, 85.20) -- (104.54, 79.20);

\path[draw=drawColor,line width= 0.4pt,line join=round,line cap=round] (121.87, 85.20) -- (121.87, 79.20);

\path[draw=drawColor,line width= 0.4pt,line join=round,line cap=round] (139.21, 85.20) -- (139.21, 79.20);

\path[draw=drawColor,line width= 0.4pt,line join=round,line cap=round] (156.54, 85.20) -- (156.54, 79.20);

\path[draw=drawColor,line width= 0.4pt,line join=round,line cap=round] (173.87, 85.20) -- (173.87, 79.20);

\path[draw=drawColor,line width= 0.4pt,line join=round,line cap=round] (191.21, 85.20) -- (191.21, 79.20);

\path[draw=drawColor,line width= 0.4pt,line join=round,line cap=round] (208.54, 85.20) -- (208.54, 79.20);

\path[draw=drawColor,line width= 0.4pt,line join=round,line cap=round] (225.88, 85.20) -- (225.88, 79.20);

\path[draw=drawColor,line width= 0.4pt,line join=round,line cap=round] (243.21, 85.20) -- (243.21, 79.20);

\node[text=drawColor,rotate= 90.00,anchor=base east,inner sep=0pt, outer sep=0pt, scale=  1.00] at ( 73.41, 73.20) {(Intercept)};

\node[text=drawColor,rotate= 90.00,anchor=base east,inner sep=0pt, outer sep=0pt, scale=  1.00] at ( 90.74, 73.20) {OBP};

\node[text=drawColor,rotate= 90.00,anchor=base east,inner sep=0pt, outer sep=0pt, scale=  1.00] at (108.08, 73.20) {SLG};

\node[text=drawColor,rotate= 90.00,anchor=base east,inner sep=0pt, outer sep=0pt, scale=  1.00] at (125.41, 73.20) {EXP};

\node[text=drawColor,rotate= 90.00,anchor=base east,inner sep=0pt, outer sep=0pt, scale=  1.00] at (142.75, 73.20) {PAYR};

\node[text=drawColor,rotate= 90.00,anchor=base east,inner sep=0pt, outer sep=0pt, scale=  1.00] at (160.08, 73.20) {MLAVG};

\node[text=drawColor,rotate= 90.00,anchor=base east,inner sep=0pt, outer sep=0pt, scale=  1.00] at (177.42, 73.20) {AVGcumLag1};

\node[text=drawColor,rotate= 90.00,anchor=base east,inner sep=0pt, outer sep=0pt, scale=  1.00] at (194.75, 73.20) {OBPcumLag1};

\node[text=drawColor,rotate= 90.00,anchor=base east,inner sep=0pt, outer sep=0pt, scale=  1.00] at (212.09, 73.20) {SLGcumLag1};

\node[text=drawColor,rotate= 90.00,anchor=base east,inner sep=0pt, outer sep=0pt, scale=  1.00] at (229.42, 73.20) {G};

\node[text=drawColor,rotate= 90.00,anchor=base east,inner sep=0pt, outer sep=0pt, scale=  1.00] at (246.75, 73.20) {YRINDEX};

\path[draw=drawColor,line width= 0.4pt,line join=round,line cap=round] ( 61.20, 88.85) -- ( 61.20,139.96);

\path[draw=drawColor,line width= 0.4pt,line join=round,line cap=round] ( 61.20, 88.85) -- ( 55.20, 88.85);

\path[draw=drawColor,line width= 0.4pt,line join=round,line cap=round] ( 61.20, 96.15) -- ( 55.20, 96.15);

\path[draw=drawColor,line width= 0.4pt,line join=round,line cap=round] ( 61.20,103.45) -- ( 55.20,103.45);

\path[draw=drawColor,line width= 0.4pt,line join=round,line cap=round] ( 61.20,110.75) -- ( 55.20,110.75);

\path[draw=drawColor,line width= 0.4pt,line join=round,line cap=round] ( 61.20,118.06) -- ( 55.20,118.06);

\path[draw=drawColor,line width= 0.4pt,line join=round,line cap=round] ( 61.20,125.36) -- ( 55.20,125.36);

\path[draw=drawColor,line width= 0.4pt,line join=round,line cap=round] ( 61.20,132.66) -- ( 55.20,132.66);

\path[draw=drawColor,line width= 0.4pt,line join=round,line cap=round] ( 61.20,139.96) -- ( 55.20,139.96);

\node[text=drawColor,anchor=base east,inner sep=0pt, outer sep=0pt, scale=  1.00] at ( 49.20, 85.31) {0.53};

\node[text=drawColor,anchor=base east,inner sep=0pt, outer sep=0pt, scale=  1.00] at ( 49.20, 92.61) {0.59};

\node[text=drawColor,anchor=base east,inner sep=0pt, outer sep=0pt, scale=  1.00] at ( 49.20, 99.91) {0.77};

\node[text=drawColor,anchor=base east,inner sep=0pt, outer sep=0pt, scale=  1.00] at ( 49.20,107.21) {0.78};

\node[text=drawColor,anchor=base east,inner sep=0pt, outer sep=0pt, scale=  1.00] at ( 49.20,114.51) {0.82};

\node[text=drawColor,anchor=base east,inner sep=0pt, outer sep=0pt, scale=  1.00] at ( 49.20,121.81) {0.82};

\node[text=drawColor,anchor=base east,inner sep=0pt, outer sep=0pt, scale=  1.00] at ( 49.20,129.12) {0.82};

\node[text=drawColor,anchor=base east,inner sep=0pt, outer sep=0pt, scale=  1.00] at ( 49.20,136.42) {0.82};

\path[draw=drawColor,line width= 0.4pt,line join=round,line cap=round] ( 61.20, 85.20) --
	(251.88, 85.20) --
	(251.88,143.61) --
	( 61.20,143.61) --
	( 61.20, 85.20);
\end{scope}
\end{tikzpicture}


\end{knitrout}

    \end{frame}     


   %{slide 10}
    \begin{frame}[fragile]{Use BIC to compare models}
      \fontsize{8}{8}\selectfont  
\begin{knitrout}
\definecolor{shadecolor}{rgb}{0.137, 0.137, 0.137}\begin{kframe}
\begin{alltt}
\hlkwd{plot}\hlstd{(bestsubsets,} \hlkwc{scale}\hlstd{=}\hlstr{"bic"}\hlstd{)}  \hlcom{# use BIC}
\end{alltt}
\end{kframe}
\input{c:/tmp/figures/unnamed-chunk-8-1.tikz}

\end{knitrout}

    \lc %{What are the best predictors}
    \end{frame}

    %{slide 11}
    \begin{frame}[fragile]{Generate the best candidate model}
    \fontsize{6}{6}\selectfont 

\begin{knitrout}
\definecolor{shadecolor}{rgb}{0.137, 0.137, 0.137}\begin{kframe}
\begin{alltt}
\hlstd{model} \hlkwb{<-} \hlkwd{lm}\hlstd{(AVG} \hlopt{~} \hlstd{OBP} \hlopt{+} \hlstd{SLG} \hlopt{+} \hlstd{AVGcumLag1} \hlopt{+} \hlstd{OBPcumLag1}
\hlopt{+} \hlstd{SLGcumLag1,} \hlkwc{data}\hlstd{=baseball)}

\hlkwd{summary}\hlstd{(model)}
\end{alltt}
\begin{verbatim}

Call:
lm(formula = AVG ~ OBP + SLG + AVGcumLag1 + OBPcumLag1 + SLGcumLag1, 
    data = baseball)

Residuals:
      Min        1Q    Median        3Q       Max 
-0.056014 -0.007723  0.000263  0.008180  0.040508 

Coefficients:
             Estimate Std. Error t value Pr(>|t|)    
(Intercept)  0.027871   0.002497   11.16   <2e-16 ***
OBP          0.498209   0.009089   54.81   <2e-16 ***
SLG          0.160826   0.004697   34.24   <2e-16 ***
AVGcumLag1   0.880348   0.011950   73.67   <2e-16 ***
OBPcumLag1  -0.476261   0.012105  -39.34   <2e-16 ***
SLGcumLag1  -0.171831   0.005547  -30.98   <2e-16 ***
---
Signif. codes:  0 '***' 0.001 '**' 0.01 '*' 0.05 '.' 0.1 ' ' 1

Residual standard error: 0.01213 on 4529 degrees of freedom
Multiple R-squared:  0.821,	Adjusted R-squared:  0.8208 
F-statistic:  4154 on 5 and 4529 DF,  p-value: < 2.2e-16
\end{verbatim}
\end{kframe}
\end{knitrout}

    \end{frame}


    %{slide 12}
    \begin{frame}[fragile]{Does the National League's Designated Hitter Rule matter?}
    \fontsize{8}{8}\selectfont  

\begin{knitrout}
\definecolor{shadecolor}{rgb}{0.137, 0.137, 0.137}\begin{kframe}
\begin{alltt}
\hlcom{#Add a the categorical variable LG and find out!}

\hlstd{model} \hlkwb{<-} \hlkwd{lm}\hlstd{(AVG} \hlopt{~} \hlstd{OBP} \hlopt{+} \hlstd{SLG} \hlopt{+} \hlstd{AVGcumLag1} \hlopt{+} \hlstd{OBPcumLag1}
\hlopt{+} \hlstd{SLGcumLag1} \hlopt{+} \hlstd{LG,} \hlkwc{data}\hlstd{=baseball)}
\end{alltt}
\end{kframe}
\end{knitrout}

    \end{frame}


    %{slide 13}
    \begin{frame}[fragile]{Does the National League's Designated Hitter Rule Matter?}
    \fontsize{8}{8}\selectfont  
    
\begin{knitrout}
\definecolor{shadecolor}{rgb}{0.137, 0.137, 0.137}\begin{kframe}
\begin{alltt}
\hlcom{# Find the rows in baseball where LG is not either NL or AL}
\hlcom{# and remove them so we can focus on the difference}
\hlcom{# between NL and AL}

\hlstd{base1} \hlkwb{<-} \hlstd{baseball[baseball}\hlopt{$}\hlstd{LG} \hlopt{==} \hlstr{"NL"} \hlopt{|} \hlstd{baseball}\hlopt{$}\hlstd{LG} \hlopt{==} \hlstr{"AL"}\hlstd{,]}

\hlstd{modelLG} \hlkwb{<-} \hlkwd{lm}\hlstd{(AVG} \hlopt{~} \hlstd{OBP} \hlopt{+} \hlstd{SLG} \hlopt{+} \hlstd{AVGcumLag1} \hlopt{+} \hlstd{OBPcumLag1}
\hlopt{+} \hlstd{SLGcumLag1} \hlopt{+} \hlstd{LG,} \hlkwc{data}\hlstd{=base1)}

\hlcom{# Look at the summary, LG is not statistically significant}
\end{alltt}
\end{kframe}
\end{knitrout}

      \lc %{What is the p-value of league?}
      \lc %{Does the designated hitter rule predict}
    \end{frame}


    %{slide 14}
    \begin{frame}[fragile]{Check for linear relationships}
    \fontsize{8}{8}\selectfont  

\begin{knitrout}
\definecolor{shadecolor}{rgb}{0.137, 0.137, 0.137}\begin{kframe}
\begin{alltt}
\hlcom{# Depending on your computer, this command may run slowly}
\hlcom{#pairs(~ AVG + OBP + SLG + AVGcumLag1 + OBPcumLag1 + SLGcumLag1, data=baseball)}
\end{alltt}
\end{kframe}
\end{knitrout}

      \begin{center}
        \includegraphics[width=3.4in]{PairsPLot.png} \\
      \end{center}

    \end{frame}


    %{slide 15}
    \begin{frame}[fragile]{Is this model really useful?}
      \begin{itemize}[<+->]
        \item Automated regression model selection methods cannot make something out of nothing. 
        \item If you omit some important variables or fail to use data transformations when they are needed, or if the assumption of linear or linearizable relationships is simply wrong, the model is a bad one, no matter what the $R^2$.  
        \item Use your own judgment and intuition about your data to try to fine-tune whatever the computer comes up with.

        \end{itemize} 
    \end{frame}


   %{slide 16}
    \begin{frame}[fragile]{A challenge}
    \fontsize{8}{8}\selectfont  

\begin{knitrout}
\definecolor{shadecolor}{rgb}{0.137, 0.137, 0.137}\begin{kframe}
\begin{alltt}
\hlcom{# Load the dataset challenge and run the following regression}

\hlstd{model} \hlkwb{<-} \hlkwd{lm}\hlstd{(y} \hlopt{~} \hlstd{x1} \hlopt{+} \hlstd{x2} \hlopt{+} \hlstd{x3} \hlopt{+} \hlstd{x4} \hlopt{+} \hlstd{x5} \hlopt{+} \hlstd{x6} \hlopt{+} \hlstd{x7}
     \hlopt{+} \hlstd{x8} \hlopt{+} \hlstd{x9} \hlopt{+} \hlstd{x10} \hlopt{+} \hlstd{x11} \hlopt{+} \hlstd{x12,} \hlkwc{data}\hlstd{=challenge)}

    \hlcom{# Generate a summary and examine R^2}
    \hlcom{# 0.21 of the variance is explained    }
\end{alltt}
\end{kframe}
\end{knitrout}

    \end{frame}


  %{slide 17}
    \begin{frame}[fragile]{Surprise!}
    \fontsize{8}{8}\selectfont  

\begin{knitrout}
\definecolor{shadecolor}{rgb}{0.137, 0.137, 0.137}\begin{kframe}
\begin{alltt}
\hlcom{# I created this data with a random number generator}
\hlcom{# You may have to run it a couple of times to get significance}

 \hlstd{y} \hlkwb{<-} \hlkwd{rnorm}\hlstd{(}\hlnum{100}\hlstd{)}
 \hlstd{x1} \hlkwb{<-} \hlkwd{rnorm}\hlstd{(}\hlnum{100}\hlstd{)}
 \hlstd{x2} \hlkwb{<-} \hlkwd{rnorm}\hlstd{(}\hlnum{100}\hlstd{)}
 \hlstd{x3} \hlkwb{<-} \hlkwd{rnorm}\hlstd{(}\hlnum{100}\hlstd{)}
 \hlstd{x4} \hlkwb{<-} \hlkwd{rnorm}\hlstd{(}\hlnum{100}\hlstd{)}
 \hlstd{x5} \hlkwb{<-} \hlkwd{rnorm}\hlstd{(}\hlnum{100}\hlstd{)}
 \hlstd{x6} \hlkwb{<-} \hlkwd{rnorm}\hlstd{(}\hlnum{100}\hlstd{)}
 \hlstd{x7} \hlkwb{<-} \hlkwd{rnorm}\hlstd{(}\hlnum{100}\hlstd{)}
 \hlstd{x8} \hlkwb{<-} \hlkwd{rnorm}\hlstd{(}\hlnum{100}\hlstd{)}
 \hlstd{x9} \hlkwb{<-} \hlkwd{rnorm}\hlstd{(}\hlnum{100}\hlstd{)}
 \hlstd{x10} \hlkwb{<-} \hlkwd{rnorm}\hlstd{(}\hlnum{100}\hlstd{)}
 \hlstd{x11} \hlkwb{<-} \hlkwd{rnorm}\hlstd{(}\hlnum{100}\hlstd{)}
 \hlstd{x12} \hlkwb{<-} \hlkwd{rnorm}\hlstd{(}\hlnum{100}\hlstd{)}
 \hlkwd{summary}\hlstd{(}\hlkwd{lm}\hlstd{(y} \hlopt{~} \hlstd{x1} \hlopt{+} \hlstd{x2} \hlopt{+} \hlstd{x3} \hlopt{+} \hlstd{x4} \hlopt{+} \hlstd{x5} \hlopt{+} \hlstd{x6} \hlopt{+} \hlstd{x7}
 \hlopt{+} \hlstd{x8} \hlopt{+} \hlstd{x9} \hlopt{+} \hlstd{x10} \hlopt{+} \hlstd{x11} \hlopt{+} \hlstd{x12))}
\end{alltt}
\begin{verbatim}

Call:
lm(formula = y ~ x1 + x2 + x3 + x4 + x5 + x6 + x7 + x8 + x9 + 
    x10 + x11 + x12)

Residuals:
     Min       1Q   Median       3Q      Max 
-2.06821 -0.58499  0.03505  0.60516  2.13938 

Coefficients:
             Estimate Std. Error t value Pr(>|t|)  
(Intercept)  0.162527   0.096608   1.682   0.0961 .
x1           0.092399   0.090793   1.018   0.3116  
x2           0.063201   0.107415   0.588   0.5578  
x3          -0.060678   0.098623  -0.615   0.5400  
x4          -0.083326   0.097662  -0.853   0.3959  
x5           0.025192   0.107542   0.234   0.8153  
x6           0.043019   0.098163   0.438   0.6623  
x7           0.138090   0.119123   1.159   0.2495  
x8           0.131719   0.108639   1.212   0.2286  
x9           0.005395   0.105470   0.051   0.9593  
x10          0.198551   0.110800   1.792   0.0766 .
x11         -0.128420   0.105479  -1.217   0.2267  
x12         -0.015154   0.108004  -0.140   0.8887  
---
Signif. codes:  0 '***' 0.001 '**' 0.01 '*' 0.05 '.' 0.1 ' ' 1

Residual standard error: 0.9543 on 87 degrees of freedom
Multiple R-squared:  0.07369,	Adjusted R-squared:  -0.05408 
F-statistic: 0.5767 on 12 and 87 DF,  p-value: 0.8553
\end{verbatim}
\begin{alltt}
    \hlcom{# 0.21 of the variance is explained by random numbers!   }
\end{alltt}
\end{kframe}
\end{knitrout}

     \lc %{What did I do wrong?}
    \end{frame}


    %{slide 18}
    \begin{frame}[fragile]{Be careful of spurious correlations and overfitting!}
      \begin{itemize}[<+->]
        \item If you have more than 1 predictor for 10-15 y values, you are likely to see spurious correlations.
        \item If you fit models with meaningless variables, you are fitting noise and will end up with an overfitted model that is not predictive going forward. 
        \item You could even end up int he American Statistical Association's Hall of Shame!
        \end{itemize} 
    \end{frame}


    %{slide 19}
    \begin{frame}{www.tylervigen.com/spurious-correlations}
      \fontsize{10}{10}\selectfont
      \begin{center}
        \includegraphics[width=3.4in]{SociologyDoctorates.png} \\
      \end{center}

      \begin{itemize}
        \item Don't fall for these!
        \item Have a great spring break!
      \end{itemize}
      
    \end{frame}

  \end{darkframes}
\end{document}
